\subsubsection{General Requirement}

Code shall not produce:

\begin{itemize}
	\item out-of-bounds accesses;
	\item use-after-free;
	\item double-free;
	\item invalid frees;
	\item null-pointer dereferences;
	\item uninitialized reads;
	\item memory leaks;
	\item access through a dangling pointer;
	\item writes through a read-only object;
	\item integer-overflow-derived allocation sizes.
\end{itemize}

For maximum saftey general rules include:

\begin{itemize}
    \item no heap allocation in core logic;
    \item no heap allocation unless needed;
    \item caller-provided buffers where possible;
    \item explicit buffer lengths;
    \item no hidden ownership transfer;
    \item no ownership implied only by naming;
    \item cleanup paths tested.
\end{itemize}

\subsubsection{Allocation Policy}

Static allocation and caller-provided storage are preferred.

Dynamic allocation shall be avoided in ordinary modules and is prohibited after initialization in deterministic or safety-critical modules.

Every allocation shall have:

\begin{itemize}
	\item one owner;
	\item one documented release path;
	\item a checked failure path;
	\item a defined maximum size;
	\item a test for boundary conditions.
\end{itemize}

\subsubsection{Allocation Arithmetic}

Allocation-size arithmetic shall be checked before allocation.

\begin{lstlisting}[language=C]
if (count > SIZE_MAX / sizeof(*items)) {
	return ERROR_SIZE_OVERFLOW;
}

items = malloc(count * sizeof(*items));
if (items == NULL) {
	return ERROR_OUT_OF_MEMORY;
}
\end{lstlisting}

\subsubsection{Ownership}

Functions shall document whether pointer parameters are:

\begin{itemize}
	\item borrowed and not retained;
	\item owned by the caller;
	\item transferred to the callee;
	\item returned with ownership transferred to the caller.
\end{itemize}
